% Part: applied-modal-logics
% Chapter: epistemic-logic
% Section: truth-at-w

\documentclass[../../../include/open-logic-section]{subfiles}

\begin{document}

\olfileid{aml}{el}{trw}

\olsection{Kabeneran ing Sawijining Jagad}

Kaya ing logika modal lumrah, saben model epistemik nemtokake !!{formula}s
endi kang bener ing saben jagad ing model iku. Kita nggunakake notasi kang
padha, ``model $\mModel{M}$ ndadekake !!{formula}~$!A$ bener ing
jagad~$w$,'' kanggo gagasan dhasar semantik relasional. Relasi kabeneran
iki ditetepake kanthi induksi lan padha karo ing logika modal lumrah
kanggo kabeh operator nonmodal.

\begin{defn}\ollabel{defn:mmodels}
  \emph{Kabeneran !!a{formula}~$!A$ ing~$w$} sajroning~$\mModel M$,
  kang dinotasikake $\mSat{M}{!A}[w]$, ditetepake kanthi induksi
  kaya mangkene:
  \begin{enumerate}
  \tagitem{prvFalse}{\indcase{!A}{\lfalse}{$\mSat{M}{\lfalse}[w]$ ora nate kelakon}.}{}
  \tagitem{prvTrue}{\indcase{!A}{\ltrue}{$\mSat{M}{\ltrue}[w]$ tansah kelakon}.}{}
  \item $\mSat{M}{p}[w]$ yen lan mung yen $w \in V(p)$
  \tagitem{prvNot}{\indcase{!A}{\lnot !B}{$\mSat{M}{\indfrm}[w]$ yen lan mung yen
    $\mSat{M}{!B}[w]$ ora kelakon}.}{}
  \tagitem{prvAnd}{\indcase{!A}{(!B \land !C)}{$\mSat{M}{\indfrm}[w]$ yen lan mung yen
    $\mSat{M}{!B}[w]$ lan $\mSat{M}{!C}[w]$}.}{}
  \tagitem{prvOr}{\indcase{!A}{(!B \lor !C)}{$\mSat{M}{\indfrm}[w]$ yen lan mung yen
    $\mSat{M}{!B}[w]$ utawa $\mSat{M}{!C}[w]$} (utawa loro-lorone).}{}
  \tagitem{prvIf}{\indcase{!A}{(!B \lif !C)}{$\mSat{M}{\indfrm}[w]$ yen lan mung yen
    $\mSat{M}{!B}[w]$ ora kelakon utawa $\mSat{M}{!C}[w]$}.}{}
  \tagitem{prvIff}{\indcase{!A}{(!B \liff !C)}{$\mSat{M}{\indfrm}[w]$ yen lan mung yen
    $\mSat{M}{!B}[w]$ lan $\mSat{M}{!C}[w]$ loro-lorone kelakon,
    utawa $\mSat{M}{!B}[w]$ lan $\mSat{M}{!C}[w]$ loro-lorone ora kelakon}.}{}
  \item\ollabel{defn:sub:mmodels-box}
    \indcase{!A}{\Knows_a !B}{$\mSat{M}{\indfrm}[w]$ yen lan mung yen
    $\mSat{M}{!B}[w']$ kanggo saben $w' \in W$ kang nyukupi $R_a ww'$}
  \end{enumerate} 
\end{defn}

Nanging, ing kene kita kudu mikirake watesan tumrap relasi
aksesibilitas. Miturut klausa~\olref{defn:sub:mmodels-box},
!!a{formula}~$\Knows_a !B$ bener ing~$w$ yen ora ana~$w'$ kang
nyukupi $R_a ww'$. Iki padha karo klausa ing logika modal lumrah:
yen sawijining jagad ora duwe penerus, kabeh rumus mawa
$\Box$ bener kanthi kosong ing jagad iku. Tafsiran iki angel
ditampa yen $\Knows$ dianggep nglambangake \emph{kawruh}.
Kita lumrahe nganggep \emph{ora ana} kahanan kang ndadekake sawijining agen
ngerti $!A$ lan $\lnot !A$ bebarengan.

Salah siji cara kanggo ngatasi prakara iki yaiku mesthekake yen
relasi aksesibilitas ing logika epistemik tansah \emph{refleksif}.
Iki kurang luwih cocog karo gagasan yen jagad nyata selaras karo
informasi kang diduweni sawijining agen. Logika epistemik biasane
nggunakake S5, nanging ana uga kang nggunakake sistem luwih ringkih,
gumantung marang apa kang dikarepake supaya diwakili dening relasi
$\Knows_a$.

\begin{figure}
  \begin{center}
    \begin{tikzpicture}[modal]
      \node[world] (w1) [label={[align=right]right:\mTrue{p}\\ \mFalse{q}}]{$w_1$}; 
      \node[world] (w2) [label={[align=right]right:\mTrue{p}\\ \mTrue{q}},
        above right=of w1]{$w_2$}; 
      \node[world] (w3) [label={[align=right]right:\mFalse{p}\\ \mFalse{q}},
        below right=of w1] {$w_3$};
      \draw[<->] (w1) to node {$a$} (w2);
      \draw[->,loop left] (w1) to node {$a,b$} (w1);
      \draw[<->] (w1) to [swap] node {$b$} (w3);
      \draw[->,loop above] (w2) to node {$a,b$} (w2) ;
      \draw[->,loop below] (w3) to node {$a,b$} (w3) ;
    \end{tikzpicture}
  \end{center}
  \caption{Sawijining model epistemik prasaja.}
  \ollabel{fig:simple}
\end{figure}

\begin{prob}
  Temtokna endi saka pratelan ing ngisor iki kang bener ing
  \olref[aml][el][trw]{fig:simple}:
  \begin{enumerate}
    \item $\mSat{M}{\lnot q}[w_1]$;
    \item $\mSat{M}{\Knows_a \lnot q}[w_1]$;
    \item $\mSat{M}{\Knows_b \lnot q}[w_1]$;
    \item $\mSat{M}{\Knows_b q \lor \Knows_b \lnot q}[w_2]$;
    \item $\mSat{M}{\Knows_a( \Knows_b q \lor \Knows_b \lnot q)}[w_2]$;
    \item $\mSat{M}{\EKnows_{\{a,b\}} \lnot q}[w_3]$;
    \end{enumerate}
\end{prob}

Sawise menehi definisi dhasar kabeneran ing sawijining jagad, kita
bisa ngetrapake pangertèn semantik liyane saka logika modal lumrah,
kayata validitas modal lan konsekuensi logis, marang cara anyar iki
kanggo napsirake operator modal.

Saiki kita uga bisa menehi syarat kabeneran kanggo operator kawruh
umum~$\CKnows_G$. Elinga saka \olref[sfr][rel][ops]{sec} yen
\emph{panutupan transitif}~$R^+$ saka relasi~$R$ ditetepake minangka
\begin{align*}
  R^+ &= \bigcup_{ n \in \mathbb{N}} R^n, 
\intertext{ing kene}
  R^0 & = R \text{ lan}\\ 
  R^{n+1} & = \Setabs{\tuple{x, z}}{\exists y (R^n xy \land Ryz)}.
\end{align*}
Banjur, yen $G$ iku klompok agen, kita netepake $R_G = ( \bigcup_{b
\in G} R_b )^+$ minangka panutupan transitif saka gabungan relasi
aksesibilitas kabeh agen ing klompok iku.

\begin{defn}
Yen $G' \subseteq G$, kita netepake $\mSat{M}{\CKnows_{G'} !A}[ w]$
yen lan mung yen kanggo saben $w'$ kang nyukupi $R_{G'} w w'$,
$\mSat{M}{!A}[w']$ kelakon.
\end{defn}

\end{document}
